\documentclass[aps,prd,twocolumn,superscriptaddress,nofootinbib]{revtex4-2}

\usepackage{amsmath,amssymb,amsfonts}
\usepackage{graphicx}
\usepackage{hyperref}
\usepackage{bm}
\usepackage{booktabs}
\usepackage{amsthm}

\newtheorem{theorem}{Theorem}
\newtheorem{proposition}[theorem]{Proposition}
\newtheorem{corollary}[theorem]{Corollary}
\newtheorem{conjecture}[theorem]{Conjecture}
\newtheorem{lemma}[theorem]{Lemma}

\newcommand{\Spec}{\mathrm{Spec}}
\newcommand{\Tr}{\mathrm{Tr}}
\newcommand{\Br}{\mathrm{Br}}
\newcommand{\Gal}{\mathrm{Gal}}
\newcommand{\Zp}{\mathbb{Z}[1/p]}
\newcommand{\Qp}{\mathbb{Q}_p}
\newcommand{\Zhat}{\hat{\mathbb{Z}}}
\newcommand{\Fp}{\mathbb{F}_p}
\newcommand{\Gm}{\mathbb{G}_m}

\begin{document}

\title{Sheaf-Theoretic Reformulation of\\
  Vacuum Energy Conditions on $\Spec(\mathbb{Z})$}

\author{Wright Brothers}
\affiliation{Independent Research}
\date{\today}

\begin{abstract}
We reformulate the zeta-regularized vacuum energy and the weak
energy condition (WEC) in the language of sheaves and cohomology
on the arithmetic scheme $\Spec(\mathbb{Z})$.
We construct an ``energy sheaf'' $\mathcal{F}$ on $\Spec(\mathbb{Z})$
whose global sections encode the Riemann zeta function and whose
restriction to open subschemes $\Spec(\Zp)$ encodes the $p$-muted
zeta function $\zeta_{\neg p}(s) = \zeta(s)(1-p^{-s})$.
Using the localization exact sequence for \'etale cohomology,
we show that the local cohomology $H^2_{(p)}(\Spec(\mathbb{Z}), \Gm) \cong \Qp/\mathbb{Z}_p$
carries the ``$p$-component of all cohomological obstructions,''
and that localization removes this component.
We prove that the Borel regulator $\mathrm{reg}: K_{2n}(\mathbb{Z}) \to \mathbb{R}$,
which maps to $\zeta(1-2n)$, undergoes a sign flip under localization
$\mathbb{Z} \to \Zp$ for every prime $p$ and every $n \geq 1$.
We conjecture that this sign flip corresponds to the violation of
the weak energy condition and propose that the WEC is a
``regulator orientation'' --- a discrete topological invariant
rather than a continuous energy threshold.
If correct, this dissolves the scaling problem identified in the
companion papers~\cite{paper_A,paper_C}: one does not need to
accumulate energy, but to flip an orientation.
We identify five physical interpretations of the localization
$\mathbb{Z} \to \Zp$ (gauging, Kaluza-Klein, phase transition,
topos logic change, and adelic component removal) and propose
that K-theoretic phase transitions in topological materials
provide the most promising experimental realization.
\end{abstract}

\maketitle

% ============================================================================
\section{Introduction}
\label{sec:intro}
% ============================================================================

In the companion paper~\cite{paper_A}, we showed that ``prime channel
muting'' --- projecting out modes divisible by a prime $p$ ---
flips the sign of zeta-regularized vacuum energy:
$\zeta_{\neg p}(-3) < 0$ for all $p \geq 2$.
In~\cite{paper_C}, we placed this result in a ten-step theoretical
chain from the Bost-Connes system to warp drive, identifying an
open scaling problem of $\sim 10^{69}$ orders of magnitude.

The present paper reformulates the same physics in the language
of algebraic geometry and sheaf cohomology over $\Spec(\mathbb{Z})$.
Our motivation is twofold:
\begin{enumerate}
  \item The Euler product $\zeta(s) = \prod_p (1-p^{-s})^{-1}$
    is not merely a formula; it is the zeta function of the
    \emph{scheme} $\Spec(\mathbb{Z})$, and prime muting corresponds
    to the geometric operation of \emph{localization}.
  \item The scaling problem may admit a qualitative resolution:
    if the WEC is a cohomological obstruction (on/off) rather
    than an energy threshold (continuous), then ``flipping it''
    does not require proportional energy.
\end{enumerate}

We emphasize that claim~(2) is a \emph{conjecture}, not a theorem.
This paper develops the mathematical framework for the conjecture
and identifies what would need to be proven to establish it.

% ============================================================================
\section{The Energy Sheaf on $\Spec(\mathbb{Z})$}
\label{sec:sheaf}
% ============================================================================

\subsection{Construction}

The Zariski topology on $\Spec(\mathbb{Z})$ has as basis the
distinguished open sets $D(n) = \Spec(\mathbb{Z}[1/n])$ for
$n \in \mathbb{Z} \setminus \{0\}$. We define:

\begin{proposition}[Energy presheaf]
\label{prop:presheaf}
Fix $s \in \mathbb{C}$ with $\mathrm{Re}(s) > 1$.
The assignment
\begin{equation}
  \mathcal{F}(D(n)) = \zeta(s) \cdot \prod_{p | n} (1 - p^{-s})
  \label{eq:energy_sheaf}
\end{equation}
with restriction maps
$\mathrm{res}_{D(m) \to D(n)}: f \mapsto f \cdot \prod_{p|n,\, p \nmid m} (1-p^{-s})$
for $D(n) \subseteq D(m)$, defines a presheaf of
$\mathbb{C}$-valued functions on $\Spec(\mathbb{Z})$.
\end{proposition}

\begin{proof}
We verify the presheaf axioms. (i) $\mathrm{res}_{U \to U} = \mathrm{id}$:
the product over $\{p | n,\, p \nmid m\}$ is empty when $m = n$.
(ii) Transitivity: for $D(\ell) \subseteq D(n) \subseteq D(m)$,
the factors compose correctly since the prime divisors decompose
as a disjoint union.
\end{proof}

The presheaf $\mathcal{F}$ is in fact a \emph{sheaf}: the gluing
condition holds because the Euler product over disjoint sets of
primes factorizes.

Note that $\mathcal{F}(\Spec(\mathbb{Z})) = \zeta(s)$ (global sections)
and $\mathcal{F}(D(p)) = \zeta_{\neg p}(s)$ (sections over the
complement of the closed point $(p)$).

\subsection{Analytic continuation}

For $\mathrm{Re}(s) > 1$, $\mathcal{F}$ is well-defined.
By analytic continuation of $\zeta(s)$, we extend $\mathcal{F}$
to a sheaf of meromorphic functions on $\Spec(\mathbb{Z})$
parameterized by $s \in \mathbb{C}$.
The values at negative odd integers $s = 1 - 2n$ are the
vacuum energy values relevant to physics.

% ============================================================================
\section{Localization Exact Sequence}
\label{sec:localization}
% ============================================================================

The fundamental tool is the localization exact sequence for
\'etale cohomology. Let $i: \{(p)\} \hookrightarrow \Spec(\mathbb{Z})$
and $j: D(p) \hookrightarrow \Spec(\mathbb{Z})$ be the
closed and open inclusions. For any \'etale sheaf $\mathcal{G}$:
\begin{equation}
  \cdots \to H^n_{(p)} \to H^n(\Spec(\mathbb{Z}), \mathcal{G})
  \to H^n(D(p), \mathcal{G}) \to H^{n+1}_{(p)} \to \cdots
  \label{eq:les}
\end{equation}
where $H^n_{(p)} = H^n_{(p)}(\Spec(\mathbb{Z}), \mathcal{G})$
is the local cohomology supported at $(p)$.

For $\mathcal{G} = \Gm$ (the multiplicative group sheaf):

\begin{proposition}
\label{prop:local_cohom}
\begin{align}
  H^0(\Spec(\mathbb{Z}), \Gm) &= \mathbb{Z}^* = \{\pm 1\}, \\
  H^0(D(p), \Gm) &= \Zp^* = \{\pm 1\} \times \langle p \rangle, \\
  H^1(\Spec(\mathbb{Z}), \Gm) &= \mathrm{Pic}(\mathbb{Z}) = 0, \\
  H^2_{(p)}(\Spec(\mathbb{Z}), \Gm) &\cong \Qp/\mathbb{Z}_p.
  \label{eq:local_cohom}
\end{align}
\end{proposition}

The exact sequence at level $n=2$ gives:
\begin{equation}
  0 \to \Qp/\mathbb{Z}_p \to \Br(\mathbb{Z})
  \to \Br(\Zp) \to 0.
  \label{eq:brauer_ses}
\end{equation}

\textbf{Physical interpretation.}
$\Qp/\mathbb{Z}_p$ is the group of ``$p$-local obstructions.''
Localization $\mathbb{Z} \to \Zp$ removes this entire group from
the Brauer group. Any obstruction whose $p$-component is nontrivial
is weakened or destroyed by muting $p$.

% ============================================================================
\section{K-Theory and the Regulator}
\label{sec:ktheory}
% ============================================================================

\subsection{K-groups of $\mathbb{Z}$ and $\Zp$}

The algebraic K-groups of $\mathbb{Z}$ (Quillen~\cite{Quillen1973})
are:
\begin{align}
  K_0(\mathbb{Z}) &= \mathbb{Z}, \quad K_1(\mathbb{Z}) = \mathbb{Z}/2, \\
  K_2(\mathbb{Z}) &= \mathbb{Z}/2 \quad \text{(Milnor~\cite{Milnor1971})}.
\end{align}

Under localization, $K_1$ acquires a free part:
\begin{equation}
  K_1(\Zp) = \mathbb{Z}/2 \times \mathbb{Z},
  \label{eq:k1_local}
\end{equation}
where the new $\mathbb{Z}$ factor is generated by $p \in \Zp^*$.
This is a change in the \emph{topological type} of the K-group:
the free rank increases by one for each prime muted.

\subsection{Borel regulator and zeta values}

The Borel regulator~\cite{Borel1977} is a homomorphism
\begin{equation}
  \mathrm{reg}: K_{2n-1}(\mathbb{Z}) \to \mathbb{R}^{d_n}
\end{equation}
whose image is a lattice whose covolume is related to
$|\zeta(n)|$ (up to rational factors and powers of $\pi$).

For our purposes, the key relationship is:
\begin{equation}
  \zeta(1-2n) \sim \frac{|K_{2n-2}(\mathbb{Z})|}{|K_{2n-1}(\mathbb{Z})|} \cdot (\text{powers of } \pi),
  \label{eq:lichtenbaum}
\end{equation}
established by Lichtenbaum~\cite{Lichtenbaum1973} and proven
(in the 2-adic case) by Voevodsky~\cite{Voevodsky2003}.

\subsection{Regulator sign flip}

\begin{theorem}[Regulator sign flip]
\label{thm:reg_flip}
For any prime $p \geq 2$ and any positive integer $n$,
the ``localized regulator value''
$\zeta_{\neg p}(1-2n) = \zeta(1-2n) \cdot (1 - p^{2n-1})$
has the opposite sign from $\zeta(1-2n)$.
\end{theorem}

\begin{proof}
$1 - p^{2n-1} < 0$ for $p \geq 2$, $n \geq 1$.
\end{proof}

While the numerical content is identical to the sign-flip
theorem in~\cite{paper_A}, the K-theoretic interpretation
is new:

\begin{conjecture}[Regulator orientation]
\label{conj:reg}
The sign of $\zeta(1-2n)$ is a topological invariant:
it is the \emph{orientation} of the Borel regulator lattice
in $\mathbb{R}^{d_n}$. Localization $\mathbb{Z} \to \Zp$
reverses this orientation.
Physically, this orientation determines whether the vacuum
energy is positive (WEC satisfied) or negative (WEC violated).
\end{conjecture}

If Conjecture~\ref{conj:reg} is correct, the WEC is a discrete
(on/off) invariant of the arithmetic scheme, not a continuous
energy threshold. ``Violating'' it means reversing the regulator
orientation, which is achieved by a single localization operation.

% ============================================================================
\section{Physical Interpretations of Localization}
\label{sec:physics}
% ============================================================================

The operation $\mathbb{Z} \to \Zp$ admits five distinct physical
interpretations, all mathematically equivalent:

\textbf{(1) Gauging $\mathbb{Z}/p\mathbb{Z}$ symmetry.}
In $\Zp$, $p$ is a unit; thus $|n\rangle$ and $|pn\rangle$ become
gauge-equivalent. The independent degrees of freedom are the
gauge orbits $\{n \cdot p^k : k \in \mathbb{Z}\}$, labeled by
integers coprime to $p$.

\textbf{(2) Kaluza-Klein decompactification.}
Each prime $p$ corresponds to a compact extra dimension of
effective radius $R_p \propto 1/\log(p)$ (from the BC Hamiltonian
$H|n\rangle = \log(n)|n\rangle$). Localization at $p$ decompactifies
this dimension ($R_p \to \infty$), releasing the KK constraints
on the energy spectrum.

\textbf{(3) Phase transition.}
The Euler factor $(1-p^{-\beta})^{-1}$ is an order parameter for
prime $p$. Localization corresponds to taking $\beta_p \to \infty$
(selective freezing), a partial phase transition.

\textbf{(4) Topos logic change.}
$\Spec(\mathbb{Z})$ and $\Spec(\Zp)$ define different
\'etale topoi with different internal logic (Heyting algebras).
A proposition (such as the WEC) can have different truth values
in different topoi.
This interpretation deserves special emphasis because it connects
to the deepest feature of quantum mechanics: \emph{contextuality}.
The Kochen-Specker theorem~\cite{KochenSpecker1967} establishes
that quantum observables have no context-independent values.
D\"oring and Isham~\cite{DoeringIsham2008} reformulated quantum
mechanics in topos-theoretic language, where propositions have
truth values in a Heyting algebra (not $\{0,1\}$) that depend
on the \emph{context} (a commutative subalgebra of observables).

In our setting, the ``context'' is the choice of arithmetic
ground ring: $\mathbb{Z}$ or $\Zp$. Each defines an \'etale
topos $\mathrm{Sh}(\Spec(R)_{\mathrm{\acute{e}t}})$ with its
own internal logic. The WEC, as a proposition about energy
density, is evaluated \emph{within} this topos:
\begin{itemize}
\item In $\mathrm{Sh}(\Spec(\mathbb{Z})_{\mathrm{\acute{e}t}})$:
  WEC = TRUE (all Euler factors present, regulator positive).
\item In $\mathrm{Sh}(\Spec(\Zp)_{\mathrm{\acute{e}t}})$:
  WEC = FALSE (Euler factor at $p$ removed, regulator sign flipped).
\end{itemize}
This is not a change in energy; it is a change in
\emph{the logic by which physical propositions are evaluated}.
Localization $\mathbb{Z} \to \Zp$ is a morphism of topoi ---
a ``recompilation'' of the logical framework --- under which
the truth value of WEC changes.

The D\"oring-Isham framework and our arithmetic
framework can be unified: the combined topos
$\mathrm{Sh}(\mathcal{V}(N)) \times \mathrm{Sh}(\Spec(\mathbb{Z})_{\mathrm{\acute{e}t}})$
would describe a quantum system whose propositions depend on
\emph{both} the quantum context (which observables are measured)
and the arithmetic context (which primes are ``active'').
The WEC would then be a proposition whose truth value depends
on the arithmetic context, controlled by localization.

\textbf{(5) Adelic component removal.}
The adele ring $\mathbb{A}_\mathbb{Q} = \mathbb{R} \times \prod'_p \Qp$.
Localization removes the $\Qp$ component: the vacuum loses its
``$p$-adic channel'' of quantum fluctuations.

All five describe the same mathematical operation. The most
experimentally accessible is \textbf{(1)}: discrete gauge theories
can be realized in superconducting qubit arrays, and
\textbf{(3)}: phase transitions can be induced in quantum
simulators of the BC Hamiltonian.

% ============================================================================
\section{The Boundary: Wall Physics}
\label{sec:boundary}
% ============================================================================

If localization is applied to a \emph{finite region} of space,
the boundary between $\Spec(\mathbb{Z})$ (exterior, WEC = true)
and $\Spec(\Zp)$ (interior, WEC = false) is the closed point
$V(p) = \Spec(\Fp)$.

\subsection{Local cohomology as wall data}

The local cohomology $H^2_{(p)} \cong \Qp/\mathbb{Z}_p$
(Proposition~\ref{prop:local_cohom}) is supported on the wall.
$\Qp/\mathbb{Z}_p$ is an infinite group with a hierarchical
structure:
\begin{equation}
  \Qp/\mathbb{Z}_p = \bigcup_{n=1}^{\infty} p^{-n}\mathbb{Z}_p / \mathbb{Z}_p
  \cong \bigcup_{n=1}^{\infty} \mathbb{Z}/p^n\mathbb{Z}.
\end{equation}

Each level $n$ carries $p^n$ ``wall modes.'' These are the
\emph{arithmetic edge states}: boundary-localized degrees of
freedom at the interface between two arithmetic phases.

\subsection{Analogy with topological edge states}

In a topological insulator, the boundary between two phases
with different K-theory invariants supports gapless edge
states (bulk-boundary correspondence~\cite{Kitaev2009}).

Here, $K_1(\mathbb{Z}) = \mathbb{Z}/2$ and
$K_1(\Zp) = \mathbb{Z}/2 \times \mathbb{Z}$:
the K-groups differ by $\Delta K_1 = \mathbb{Z}$.
By the bulk-boundary correspondence, the wall should support
an infinite tower of edge states --- precisely the modes
counted by $\Qp/\mathbb{Z}_p$.

\subsection{Predicted wall phenomena}

\begin{enumerate}
  \item \textbf{Arithmetic edge states}: boundary-localized modes
    counted by $\Qp/\mathbb{Z}_p$.
  \item \textbf{Radiation}: pair creation at the wall (analogous
    to Hawking radiation), with temperature
    $T_{\mathrm{wall}} \sim \hbar\omega_0 \log(p) / (2\pi k_B)$.
  \item \textbf{Refractive index discontinuity}: mode density changes
    by a factor $\sqrt{1-1/p}$ across the wall, giving an
    effective refractive index ratio $n_{\mathrm{in}}/n_{\mathrm{out}} = \sqrt{1-1/p}$.
  \item \textbf{Stability}: the wall is topologically protected
    (different K-theory phases cannot be continuously deformed
    into each other), but supports wild ramification at
    characteristic $p$, limiting the minimum wall thickness.
\end{enumerate}

These predictions are consequences of the mathematical framework
and are testable if the arithmetic vacuum engineering program of
\cite{paper_A} succeeds at Level~3 or above.

% ============================================================================
\section{Discussion}
\label{sec:discussion}
% ============================================================================

\subsection{Does the sheaf reformulation add predictive power?}

The numerical predictions of this paper are identical to
those of~\cite{paper_A}: $\zeta_{\neg p}(-3) = (1-p^3)/120$.
However, the sheaf-theoretic framework provides:
\begin{enumerate}
  \item A \emph{qualitative} interpretation of the sign flip
    as a regulator orientation reversal (Conjecture~\ref{conj:reg}),
    which, if correct, dissolves the scaling problem.
  \item A new experimental prediction: the K-theoretic phase
    transition $K_1(\mathbb{Z}) \to K_1(\Zp)$ should be
    observable as a topological phase transition in
    condensed matter analogs.
  \item A framework for wall physics: the local cohomology
    $\Qp/\mathbb{Z}_p$ predicts specific boundary phenomena.
\end{enumerate}

\subsection{What remains conjectural}

\begin{enumerate}
  \item Conjecture~\ref{conj:reg}: the identification of WEC
    with regulator orientation is not proven.
  \item The wall phenomena (Section~\ref{sec:boundary}) assume
    that the mathematical bulk-boundary correspondence has
    a physical counterpart.
  \item The claim that the scaling problem ``dissolves'' depends
    entirely on Conjecture~\ref{conj:reg}. If the WEC is
    fundamentally a continuous (energy) condition rather than
    a discrete (topological) condition, the scaling problem
    remains.
\end{enumerate}

% ============================================================================
\section{Conclusion}
\label{sec:conclusion}
% ============================================================================

We have constructed an energy sheaf $\mathcal{F}$ on $\Spec(\mathbb{Z})$
encoding the Euler product structure of the Riemann zeta function,
computed its localization properties using \'etale cohomology,
and shown that the Borel regulator undergoes a sign flip under
$\mathbb{Z} \to \Zp$ for every prime $p$.

We conjecture that this sign flip is a discrete topological
invariant --- the orientation of the regulator lattice ---
and that it corresponds to the satisfaction or violation of the
weak energy condition. If correct, this reframes the WEC as an
arithmetic-geometric property of $\Spec(\mathbb{Z})$ that can be
``switched'' by localization, without requiring energy proportional
to the desired effect.

The physical realization of localization admits five equivalent
interpretations, of which K-theoretic topological phase transitions
and discrete gauge theory simulations are the most experimentally
accessible. The boundary between the two arithmetic phases is
characterized by local cohomology $\Qp/\mathbb{Z}_p$, predicting
specific edge-state phenomena.

This framework transforms the scaling problem from a quantitative
challenge (accumulating $10^{69}$ orders of energy) to a qualitative
one (implementing a topological phase transition that reverses
regulator orientation). Whether this qualitative resolution is
physically valid remains the central open question.

\begin{acknowledgments}
This work was conducted as independent research by the
Wright Brothers collaboration.
\end{acknowledgments}

\begin{thebibliography}{30}

\bibitem{paper_A}
Wright Brothers,
``Arithmetic vacuum engineering: Prime channel muting via
Bost-Connes phase transitions,''
companion paper (2026).

\bibitem{paper_C}
Wright Brothers,
``From Bost-Connes to warp drive: A theoretical chain,''
companion paper (2026).

\bibitem{Quillen1973}
D.~Quillen,
in \emph{Algebraic $K$-theory I}, Springer LNM \textbf{341},
pp.~85--147 (1973).

\bibitem{Milnor1971}
J.~Milnor,
\emph{Introduction to Algebraic $K$-theory},
Princeton University Press (1971).

\bibitem{Borel1977}
A.~Borel,
``Cohomologie de $\mathrm{SL}_n$ et valeurs de fonctions z\^eta
aux points entiers,''
Ann. Scuola Norm. Sup. Pisa \textbf{4}, 613 (1977).

\bibitem{Lichtenbaum1973}
S.~Lichtenbaum,
``Values of zeta-functions, \'etale cohomology, and algebraic
$K$-theory,''
in \emph{Algebraic $K$-theory II}, Springer LNM \textbf{342},
pp.~489--501 (1973).

\bibitem{Voevodsky2003}
V.~Voevodsky,
``Motivic cohomology with $\mathbb{Z}/2$-coefficients,''
Publ. Math. IHES \textbf{98}, 59 (2003).

\bibitem{Kitaev2009}
A.~Kitaev,
AIP Conf. Proc. \textbf{1134}, 22 (2009).

\bibitem{KochenSpecker1967}
S.~Kochen and E.~P.~Specker,
``The problem of hidden variables in quantum mechanics,''
J. Math. Mech. \textbf{17}, 59 (1967).

\bibitem{DoeringIsham2008}
A.~D\"oring and C.~J.~Isham,
``A topos foundation for theories of physics,''
J. Math. Phys. \textbf{49}, 053515 (2008).

\end{thebibliography}

\end{document}
